\documentclass[a4paper,10pt]{report}\input{../0.Préambule/Préambule_cours_prof}\begin{document}

\pagestyle{empty}
\newpage
\section*{Proportionnalité - Nombres relatifs \hspace{2cm} Sujet A}
\addcontentsline{toc}{section}{Proportionnalité - Nombres relatifs Sujet A}

\noindent \textit{L'usage de la calculatrice est \textbf{interdite}.} \hfill \textit{Durée : 55 minutes}

\medskip
\noindent \textit{Ceci est un questionnaire a choix multiples (QCM). Pour chaque question, il y a une ou plusieurs bonnes réponses.}

\medskip
\noindent \textit{Une bonne réponse rapporte un point, un mauvaise réponse enlève un tiers de point et une absence de réponse ne rapporte ni n'enlève aucun point.}

\subsection*{Nombre relatifs}
\addcontentsline{toc}{subsection}{Nombre relatifs}

\bigskip
\noindent \textbf{Affirmation 1 : } Le résultat de l'opération $(-19)+(+11)$ est :

\medskip
\begin{tabular}{p{3.8cm}p{3.8cm}p{3.8cm}p{3.8cm}}
\textbf{A. } $(-30)$&
\textbf{B. } $(+30)$&
\textbf{C. } $(-8)$&
\textbf{D. } $(-12)$\\
\end{tabular}

\bigskip
\noindent \textbf{Affirmation 2 : } $(-5)$ et $(+5)$ sont deux nombres :

\medskip
\begin{tabular}{p{3.8cm}p{3.8cm}p{3.8cm}p{3.8cm}}
\textbf{A. } contraires&
\textbf{B. } inverses&
\textbf{C. } symétriques&
\textbf{D. } opposés\\
\end{tabular}

\bigskip
\noindent \textbf{Affirmation 3 : } $-13$ est le résultat de l'opération :

\medskip
\begin{tabular}{p{3.8cm}p{3.8cm}p{3.8cm}p{3.8cm}}
\textbf{A. } $(-11)-(-2)$&
\textbf{B. } $(-2)-(-11)$&
\textbf{C. } $(+2)-(-11)$&
\textbf{D. } $(-11)-(+2)$\\
\end{tabular}

\bigskip
\noindent \textbf{Affirmation 4 : } L'expression $(-5)-(+4)$ peur s'écrire :

\medskip
\begin{tabular}{p{3.8cm}p{3.8cm}p{3.8cm}p{3.8cm}}
\textbf{A. } $5+4$&
\textbf{B. } $5-4$&
\textbf{C. } $-5+4$&
\textbf{D. } $-5-4$\\
\end{tabular}

\bigskip
\noindent \textbf{Affirmation 5 : } Le résultat du calcul $-14-4+8+7-8$ est :

\medskip
\begin{tabular}{p{3.8cm}p{3.8cm}p{3.8cm}p{3.8cm}}
\textbf{A. } $6$&
\textbf{B. } $-11$&
\textbf{C. } $-25$&
\textbf{D. } $12$\\
\end{tabular}

\bigskip
\noindent \textbf{Affirmation 6 : } Dans un QCM de dix questions, une réponse juste rapporte 3 points, une absence de réponse 0 point et une mauvaise réponse enlève 1 points. Marlène à 2 bonnes réponses et 7 mauvaises. Quelle est sa note ?

\medskip
\begin{tabular}{p{3.8cm}p{3.8cm}p{3.8cm}p{3.8cm}}
\textbf{A. } $-1$&
\textbf{B. } $-5$&
\textbf{C. } $-2$&
\textbf{D. } $6$\\
\end{tabular}

\bigskip
\noindent \textbf{Affirmation 7 : } $(+2)-(+3)+(+4)$ est égale à :

\medskip
\begin{tabular}{p{3.8cm}p{3.8cm}p{3.8cm}p{3.8cm}}
\textbf{A. } $(+2)+(+3)+(-4)$&
\textbf{B. } $(+2)+(-3)+(+4)$&
\textbf{C. } $(+2)+(-3)+(-4)$&
\textbf{D. } $(-2)+(-3)+(-4)$\\
\end{tabular}

\bigskip
\noindent \textbf{Affirmation 8 : } Une somme algébrique est une suite:

\medskip
\begin{tabular}{p{3.8cm}p{3.8cm}p{3.8cm}p{3.8cm}}
\textbf{A. } de multiplication et de division&
\textbf{B. } d'addition et de division&
\textbf{C. } de multiplication et de soustraction&
\textbf{D. } d'addition et de soustraction\\
\end{tabular}

\bigskip
\noindent Pour les affirmations $9$ et $10$, on s'intéresse au programme de calcul suivant :

\begin{center}
\renewcommand{\arraystretch}{1.3}
\begin{tabularx}{5.5cm}{|*{1}{>{\arraybackslash}X|}}
\hline
$\bullet$ \texttt{Choisis un nombre} \\
$\bullet$ \texttt{Ajoute } $-3$ \\
$\bullet$ \texttt{Retire} $-1,5$ \\
$\bullet$ \texttt{Donne l'opposé du résultat} \\\hline
\end{tabularx}
\end{center}

\bigskip
\noindent \textbf{Affirmation 9 : } Dans ce programme de calcul, si on choisit $5$ comme nombre de départ, quel nombre obtient-on a l'arrivée ?

\medskip
\begin{tabular}{p{3.8cm}p{3.8cm}p{3.8cm}p{3.8cm}}
\textbf{A. } $6,5$&
\textbf{B. } $-6,5$&
\textbf{C. } $-3,5$&
\textbf{D. } $11,5$\\
\end{tabular}

\bigskip
\noindent \textbf{Affirmation 10 : } Dans ce programme de calcul, si on choisit $-3$ comme nombre de départ, quel nombre obtient-on à l'arrivée ?:

\medskip
\begin{tabular}{p{3.8cm}p{3.8cm}p{3.8cm}p{3.8cm}}
\textbf{A. } $4,5$&
\textbf{B. } $-4,5$&
\textbf{C. } $1,5$&
\textbf{D. } $-1,5$\\
\end{tabular}

\subsection*{Proportionnalité}
\addcontentsline{toc}{subsection}{Proportionnalité}

\bigskip
\noindent \textbf{Affirmation 11 : } : Quel est le coefficient de \\proportionnalité de ce tableau de proportionnalité.

\vspace{-1cm}
\begin{flushright}
\renewcommand{\arraystretch}{1.4}
\begin{tabularx}{8cm}{|>{\columncolor{gray!50}}p{3cm}|*{2}{>{\centering \arraybackslash}X|}}
\hline
Nombre de séances & $4$ & $14$ \\\hline
Prix à payé (en \euro) & $32$ & $112$ \\\hline
\end{tabularx}
\end{flushright}

\medskip
\begin{tabular}{p{3.8cm}p{3.8cm}p{3.8cm}p{3.8cm}}
\textbf{A. } $10$&
\textbf{B. } $8$&
\textbf{C. } $80$&
\textbf{D. } on ne peut pas savoir\\
\end{tabular}

\bigskip
\noindent \textbf{Affirmation 12 : } Complète la case manquante dans le \\tableau ci-contre :

\vspace{-1cm}
\begin{flushright}
\renewcommand{\arraystretch}{1.4}
\begin{tabularx}{8cm}{|>{\columncolor{gray!50}}p{4cm}|*{2}{>{\centering \arraybackslash}X|}}
\hline
Coté du carré (en cm) & $2$ & $5$ \\\hline
Aire du carré (en cm$^2$) & $4$ &  \\\hline
\end{tabularx}
\end{flushright}

\medskip
\begin{tabular}{p{3.8cm}p{3.8cm}p{3.8cm}p{3.8cm}}
\textbf{A. } $10$&
\textbf{B. } $25$&
\textbf{C. } $4,5$&
\textbf{D. } $7$\\
\end{tabular}

\bigskip
\noindent \textbf{Affirmation 13 : } Complète la case manquante dans le \\tableau de proportionnalité ci-dessous :

\vspace{-1cm}
\begin{flushright}
\renewcommand{\arraystretch}{1.4}
\begin{tabularx}{8cm}{|>{\columncolor{gray!50}}p{4cm}|*{2}{>{\centering \arraybackslash}X|}}
\hline
Coté du carré (en cm) & $8$ & $21$ \\\hline
Périmètre du carré (en cm) & $32$ &  \\\hline
\end{tabularx}
\end{flushright}

\medskip
\begin{tabular}{p{3.8cm}p{3.8cm}p{3.8cm}p{3.8cm}}
\textbf{A. } $12$&
\textbf{B. } $45$&
\textbf{C. } $84$&
\textbf{D. } $42$\\
\end{tabular}

\bigskip
\noindent \textbf{Affirmation 14 : } Parmi les grandeurs suivantes, laquelle est proportionnelle au volume d'eau de mer. Le volume de :

\medskip
\begin{tabular}{p{3.8cm}p{3.8cm}p{3.8cm}p{3.8cm}}
\textbf{A. } sable sur la plage &
\textbf{B. } sel dans l'eau &
\textbf{C. } sel dans la soupe &
\textbf{D. } plastique sur Terre \\
\end{tabular}

\bigskip
\noindent \textbf{Affirmation 15 : } Sur les tableaux ci dessous, combien y a t-il de tableau de proportionnalité ?:

\begin{tabularx}{\linewidth}{*{3}{>{\centering \arraybackslash}X}}
\renewcommand{\arraystretch}{1.5}\begin{tabular}{|c|c|c|}
\hline
$10$ & $15$ & $30$ \\\hline
$15$ & $25$ & $50$ \\\hline
\end{tabular} &
\renewcommand{\arraystretch}{1.5}\begin{tabular}{|c|c|c|}
\hline
$20$ & $60$ & $80$ \\\hline
$50$ & $150$ & $220$ \\\hline
\end{tabular} &
\renewcommand{\arraystretch}{1.5}\begin{tabular}{|c|c|}
\hline
$8$ & $15$ \\\hline
$20$ & $40$ \\\hline
\end{tabular} \\
\end{tabularx}

\medskip
\begin{tabular}{p{3.8cm}p{3.8cm}p{3.8cm}p{3.8cm}}
\textbf{A. } $0$&
\textbf{B. } $1$&
\textbf{C. } $2$&
\textbf{D. } $3$\\
\end{tabular}

\bigskip
\noindent Pour les affirmations \textbf{16}, \textbf{17} et \textbf{18}, on s'intéresse au tableau de proportionnalité suivant :

\begin{center}
\renewcommand{\arraystretch}{1.5}
\begin{tabularx}{7.5cm}{|*{6}{>{\centering \arraybackslash}X|}}
\hline
$6$ & $9$ & $15$ & \bcbook & $30$ & \\\hline
 \bctrefle & $27$ & & $63$ & \bcplume & $84$\\\hline
\end{tabularx}
\end{center}

\bigskip
\noindent \textbf{Affirmation 16 : } La valeur de \quad \bcbook \quad peut être obtenue avec le calcul :

\medskip
\begin{tabular}{p{3.8cm}p{3.8cm}p{3.8cm}p{3.8cm}}
\textbf{A. } $\dfrac{9 \times 27}{63}$&
\textbf{B. } $\dfrac{9 \times 63}{27}$&
\textbf{C. } $\dfrac{27 \times 63}{9}$&
\textbf{D. } $\dfrac{27}{63 \times 9}$\\
\end{tabular}

\bigskip
\noindent \textbf{Affirmation 17 : } La valeur de \quad \bctrefle \quad est :

\medskip
\begin{tabular}{p{3.8cm}p{3.8cm}p{3.8cm}p{3.8cm}}
\textbf{A. } $7$&
\textbf{B. } $18$&
\textbf{C. } $2,54$&
\textbf{D. } $\dfrac{4}{3}$\\
\end{tabular}

\bigskip
\noindent \textbf{Affirmation 18 : } La valeur de \quad \bcplume \quad est :

\medskip
\begin{tabular}{p{3.8cm}p{3.8cm}p{3.8cm}p{3.8cm}}
\textbf{A. } inférieur à $30$&
\textbf{B. } supérieur à $300$&
\textbf{C. } entre $30$ et $100$&
\textbf{D. } entre $100$ et $300$\\
\end{tabular}

\bigskip
\noindent La pâtissière a pesé ses beignets et a trouvé : \quad $2$ beignets $\rightarrow 300$g et $3$ beignets $\rightarrow 450$g. 

\bigskip
\noindent \textbf{Affirmation 19 : } Combien pèsent $5$ beignets ?

\medskip
\begin{tabular}{p{3.8cm}p{3.8cm}p{3.8cm}p{3.8cm}}
\textbf{A. } $600$&
\textbf{B. } $750$&
\textbf{C. } $150$&
\textbf{D. } $900$\\
\end{tabular}

\bigskip
\noindent \textbf{Affirmation 20 : } Combien pèse $1$ beignets ?

\medskip
\begin{tabular}{p{3.8cm}p{3.8cm}p{3.8cm}p{3.8cm}}
\textbf{A. } $150$&
\textbf{B. } $100$&
\textbf{C. } $125$&
\textbf{D. } on ne peut pas savoir\\
\end{tabular}

\newpage
\section*{Proportionnalité - Nombres relatifs \hspace{2cm} Sujet B}
\addcontentsline{toc}{section}{Proportionnalité - Nombres relatifs Sujet B}

\noindent \textit{L'usage de la calculatrice est \textbf{interdite}.} \hfill \textit{Durée : 55 minutes}

\medskip
\noindent \textit{Ceci est un questionnaire a choix multiples (QCM). Pour chaque question, il y a une ou plusieurs bonnes réponses.}

\medskip
\noindent \textit{Une bonne réponse rapporte un point, un mauvaise réponse enlève un tiers de point et une absence de réponse ne rapporte ni n'enlève aucun point.}

\subsection*{Nombre relatifs}
\addcontentsline{toc}{subsection}{Nombre relatifs}

\bigskip
\noindent \textbf{Affirmation 1 : } $(-5)$ et $(+5)$ sont deux nombres :

\medskip
\begin{tabular}{p{3.8cm}p{3.8cm}p{3.8cm}p{3.8cm}}
\textbf{A. } contraires&
\textbf{B. } symétriques&
\textbf{C. } inverses&
\textbf{D. } opposés\\
\end{tabular}

\bigskip
\noindent \textbf{Affirmation 2 : } Le résultat de l'opération $(-19)+(+11)$ est :

\medskip
\begin{tabular}{p{3.8cm}p{3.8cm}p{3.8cm}p{3.8cm}}
\textbf{A. } $(-30)$&
\textbf{B. } $(-8)$&
\textbf{C. } $(+30)$&
\textbf{D. } $(-12)$\\
\end{tabular}

\bigskip
\noindent \textbf{Affirmation 3 : } L'expression $(-5)-(+4)$ peur s'écrire :

\medskip
\begin{tabular}{p{3.8cm}p{3.8cm}p{3.8cm}p{3.8cm}}
\textbf{A. } $5+4$&
\textbf{B. } $-5+4$&
\textbf{C. } $5-4$&
\textbf{D. } $-5-4$\\
\end{tabular}

\bigskip
\noindent \textbf{Affirmation 4 : } $-13$ est le résultat de l'opération :

\medskip
\begin{tabular}{p{3.8cm}p{3.8cm}p{3.8cm}p{3.8cm}}
\textbf{A. } $(-11)-(-2)$&
\textbf{B. } $(+2)-(-11)$&
\textbf{C. } $(-2)-(-11)$&
\textbf{D. } $(-11)-(+2)$\\
\end{tabular}

\bigskip
\noindent \textbf{Affirmation 5 : } Le résultat du calcul $-14-4+8+7-8$ est :

\medskip
\begin{tabular}{p{3.8cm}p{3.8cm}p{3.8cm}p{3.8cm}}
\textbf{A. } $6$&
\textbf{B. } $-25$&
\textbf{C. } $-11$&
\textbf{D. } $12$\\
\end{tabular}

\bigskip
\noindent \textbf{Affirmation 6 : } Dans un QCM de dix questions, une réponse juste rapporte 3 points, une absence de réponse 0 point et une mauvaise réponse enlève 1 points. Marlène à 2 bonnes réponses et 7 mauvaises. Quelle est sa note ?

\medskip
\begin{tabular}{p{3.8cm}p{3.8cm}p{3.8cm}p{3.8cm}}
\textbf{A. } $-1$&
\textbf{B. } $-2$&
\textbf{C. } $-5$&
\textbf{D. } $6$\\
\end{tabular}

\bigskip
\noindent \textbf{Affirmation 7 : } $(+2)-(+3)+(+4)$ est égale à :

\medskip
\begin{tabular}{p{3.8cm}p{3.8cm}p{3.8cm}p{3.8cm}}
\textbf{A. } $(+2)+(+3)+(-4)$&
\textbf{B. } $(+2)+(-3)+(-4)$&
\textbf{C. } $(+2)+(-3)+(+4)$&
\textbf{D. } $(-2)+(-3)+(-4)$\\
\end{tabular}

\bigskip
\noindent \textbf{Affirmation 8 : } Une somme algébrique est une suite:

\medskip
\begin{tabular}{p{3.8cm}p{3.8cm}p{3.8cm}p{3.8cm}}
\textbf{A. } de multiplication et de division&
\textbf{B. } de multiplication et de soustraction&
\textbf{C. } d'addition et de division&
\textbf{D. } d'addition et de soustraction\\
\end{tabular}

\bigskip
\noindent Pour les affirmations $9$ et $10$, on s'intéresse au programme de calcul suivant :

\begin{center}
\renewcommand{\arraystretch}{1.3}
\begin{tabularx}{5.5cm}{|*{1}{>{\arraybackslash}X|}}
\hline
$\bullet$ \texttt{Choisis un nombre} \\
$\bullet$ \texttt{Ajoute } $-3$ \\
$\bullet$ \texttt{Retire} $-1,5$ \\
$\bullet$ \texttt{Donne l'opposé du résultat} \\\hline
\end{tabularx}
\end{center}

\bigskip
\noindent \textbf{Affirmation 9 : } Dans ce programme de calcul, si on choisit $-3$ comme nombre de départ, quel nombre obtient-on à l'arrivée ?:

\medskip
\begin{tabular}{p{3.8cm}p{3.8cm}p{3.8cm}p{3.8cm}}
\textbf{A. } $4,5$&
\textbf{B. } $1,5$&
\textbf{C. } $-4,5$&
\textbf{D. } $-1,5$\\
\end{tabular}

\bigskip
\noindent \textbf{Affirmation 10 : } Dans ce programme de calcul, si on choisit $5$ comme nombre de départ, quel nombre obtient-on a l'arrivée ?

\medskip
\begin{tabular}{p{3.8cm}p{3.8cm}p{3.8cm}p{3.8cm}}
\textbf{A. } $6,5$&
\textbf{B. } $-3,5$&
\textbf{C. } $-6,5$&
\textbf{D. } $11,5$\\
\end{tabular}

\subsection*{Proportionnalité}
\addcontentsline{toc}{subsection}{Proportionnalité}

\bigskip
\noindent \textbf{Affirmation 11 : } Complète la case manquante dans le \\tableau ci-contre :

\vspace{-1cm}
\begin{flushright}
\renewcommand{\arraystretch}{1.4}
\begin{tabularx}{8cm}{|>{\columncolor{gray!50}}p{4cm}|*{2}{>{\centering \arraybackslash}X|}}
\hline
Coté du carré (en cm) & $2$ & $5$ \\\hline
Aire du carré (en cm$^2$) & $4$ &  \\\hline
\end{tabularx}
\end{flushright}

\medskip
\begin{tabular}{p{3.8cm}p{3.8cm}p{3.8cm}p{3.8cm}}
\textbf{A. } $10$&
\textbf{B. } $4,5$&
\textbf{C. } $25$&
\textbf{D. } $7$\\
\end{tabular}

\bigskip
\noindent \textbf{Affirmation 12 : } Complète la case manquante dans le \\tableau de proportionnalité ci-dessous :

\vspace{-1cm}
\begin{flushright}
\renewcommand{\arraystretch}{1.4}
\begin{tabularx}{8cm}{|>{\columncolor{gray!50}}p{4cm}|*{2}{>{\centering \arraybackslash}X|}}
\hline
Coté du carré (en cm) & $8$ & $21$ \\\hline
Périmètre du carré (en cm) & $32$ &  \\\hline
\end{tabularx}
\end{flushright}

\medskip
\begin{tabular}{p{3.8cm}p{3.8cm}p{3.8cm}p{3.8cm}}
\textbf{A. } $12$&
\textbf{B. } $84$&
\textbf{C. } $45$&
\textbf{D. } $42$\\
\end{tabular}

\bigskip
\noindent \textbf{Affirmation 13 : } : Quel est le coefficient de \\proportionnalité de ce tableau de proportionnalité.

\vspace{-1cm}
\begin{flushright}
\renewcommand{\arraystretch}{1.4}
\begin{tabularx}{8cm}{|>{\columncolor{gray!50}}p{3cm}|*{2}{>{\centering \arraybackslash}X|}}
\hline
Nombre de séances & $4$ & $14$ \\\hline
Prix à payé (en \euro) & $32$ & $112$ \\\hline
\end{tabularx}
\end{flushright}

\medskip
\begin{tabular}{p{3.8cm}p{3.8cm}p{3.8cm}p{3.8cm}}
\textbf{A. } $10$&
\textbf{B. } $80$&
\textbf{C. } $8$&
\textbf{D. } on ne peut pas savoir\\
\end{tabular}

\bigskip
\noindent \textbf{Affirmation 14 : } Parmi les grandeurs suivantes, laquelle est proportionnelle au volume d'eau de mer. Le volume de :

\medskip
\begin{tabular}{p{3.8cm}p{3.8cm}p{3.8cm}p{3.8cm}}
\textbf{A. } sable sur la plage &
\textbf{B. } sel dans la soupe &
\textbf{C. } sel dans l'eau &
\textbf{D. } plastique sur Terre \\
\end{tabular}

\bigskip
\noindent \textbf{Affirmation 15 : } Sur les tableaux ci dessous, combien y a t-il de tableau de proportionnalité ?:

\begin{tabularx}{\linewidth}{*{3}{>{\centering \arraybackslash}X}}
\renewcommand{\arraystretch}{1.5}\begin{tabular}{|c|c|c|}
\hline
$10$ & $15$ & $30$ \\\hline
$15$ & $25$ & $50$ \\\hline
\end{tabular} &
\renewcommand{\arraystretch}{1.5}\begin{tabular}{|c|c|c|}
\hline
$20$ & $60$ & $80$ \\\hline
$50$ & $150$ & $220$ \\\hline
\end{tabular} &
\renewcommand{\arraystretch}{1.5}\begin{tabular}{|c|c|}
\hline
$8$ & $15$ \\\hline
$20$ & $40$ \\\hline
\end{tabular} \\
\end{tabularx}

\medskip
\begin{tabular}{p{3.8cm}p{3.8cm}p{3.8cm}p{3.8cm}}
\textbf{A. } $0$&
\textbf{B. } $2$&
\textbf{C. } $1$&
\textbf{D. } $3$\\
\end{tabular}

\bigskip
\noindent Pour les affirmations \textbf{16}, \textbf{17} et \textbf{18}, on s'intéresse au tableau de proportionnalité suivant :

\begin{center}
\renewcommand{\arraystretch}{1.5}
\begin{tabularx}{7.5cm}{|*{6}{>{\centering \arraybackslash}X|}}
\hline
$6$ & $9$ & $15$ & \bctrefle& $30$ & \\\hline
 \bcbook  & $27$ & & $63$ & \bcplume & $84$\\\hline
\end{tabularx}
\end{center}


\bigskip
\noindent \textbf{Affirmation 16 : } La valeur de \quad \bcbook \quad est :

\medskip
\begin{tabular}{p{3.8cm}p{3.8cm}p{3.8cm}p{3.8cm}}
\textbf{A. } $7$&
\textbf{C. } $18$&
\textbf{B. } $2,54$&
\textbf{D. } $\dfrac{4}{3}$\\
\end{tabular}

\bigskip
\noindent \textbf{Affirmation 17 : } La valeur de \quad \bctrefle\quad peut être obtenue avec le calcul :

\medskip
\begin{tabular}{p{3.8cm}p{3.8cm}p{3.8cm}p{3.8cm}}
\textbf{A. } $\dfrac{9 \times 27}{63}$&
\textbf{B. } $\dfrac{27 \times 63}{9}$&
\textbf{C. } $\dfrac{9 \times 63}{27}$&
\textbf{D. } $\dfrac{27}{63 \times 9}$\\
\end{tabular}

\bigskip
\noindent \textbf{Affirmation 18 : } La valeur de \quad \bcplume \quad est :

\medskip
\begin{tabular}{p{3.8cm}p{3.8cm}p{3.8cm}p{3.8cm}}
\textbf{A. } inférieur à $30$&
\textbf{B. } entre $30$ et $100$&
\textbf{C. } supérieur à $300$&
\textbf{D. } entre $100$ et $300$\\
\end{tabular}

\bigskip
\noindent La pâtissière a pesé ses beignets et a trouvé : \quad $2$ beignets $\rightarrow 300$g et $3$ beignets $\rightarrow 450$g. 

\bigskip
\noindent \textbf{Affirmation 19 : } Combien pèse $1$ beignets ?

\medskip
\begin{tabular}{p{3.8cm}p{3.8cm}p{3.8cm}p{3.8cm}}
\textbf{A. } $150$&
\textbf{B. } $125$&
\textbf{C. } $100$&
\textbf{D. } on ne peut pas savoir\\
\end{tabular}

\bigskip
\noindent \textbf{Affirmation 20 : } Combien pèsent $5$ beignets ?

\medskip
\begin{tabular}{p{3.8cm}p{3.8cm}p{3.8cm}p{3.8cm}}
\textbf{A. } $600$&
\textbf{B. } $150$&
\textbf{C. } $750$&
\textbf{D. } $900$\\
\end{tabular}



\end{document}
